% ============================================================================
\chapter{Hecke 算子}
\label{ch:hecke-operators}
% ============================================================================

% ----------------------------------------------------------------------------
\section{双陪集与 Hecke 算子的定义}
% ----------------------------------------------------------------------------

\begin{definition}[Hecke 算子]
对素数 $p$，作用在 $\Mk_k(\SL_2(\Z))$ 上的 \textbf{Hecke 算子} $T_p$ 定义为
\[
  (T_p f)(\tau) = p^{k-1} f(p\tau)
  + \frac{1}{p} \sum_{j=0}^{p-1} f\!\left(\frac{\tau + j}{p}\right).
\]
等价地，用 $q$-展开：若 $f = \sum a_n q^n$，则
\[
  T_p f = \sum_{n=0}^{\infty} \bigl( a_{np} + p^{k-1} a_{n/p} \bigr) q^n,
\]
其中约定 $n/p \notin \Z$ 时 $a_{n/p} = 0$。
\end{definition}

% ----------------------------------------------------------------------------
\section{Hecke 算子的性质}
% ----------------------------------------------------------------------------

% TODO: T_m T_n = T_{mn} 当 gcd(m,n)=1，Hecke 代数的交换性

% ----------------------------------------------------------------------------
\section{Hecke 本征形式}
% ----------------------------------------------------------------------------

% TODO: 本征形式的定义、a₁ = 1 的规范化、Fourier 系数的乘性

% ----------------------------------------------------------------------------
\section{同余子群上的 Hecke 算子}
% ----------------------------------------------------------------------------

% TODO: T_n 在 Γ₀(N) 上的定义，Diamond 算子 ⟨d⟩ 与 T_p 的关系

% ----------------------------------------------------------------------------
\section{习题}
% ----------------------------------------------------------------------------
